\documentclass{article}
\usepackage[utf8]{inputenc}
\usepackage[spanish]{babel}
\usepackage{amsfonts} 
\usepackage{amsmath}
\usepackage{amsthm} 
\usepackage{amssymb}
\usepackage{nicefrac}
\usepackage{enumitem}
\usepackage{fontawesome}
\usepackage{setspace}
\usepackage{tikz-cd}
\usepackage{upgreek}
\usepackage{diffcoeff}  
\usepackage{graphicx}
\usepackage{mathtools}



\definecolor{azul}   {HTML}{1B4F72}
\definecolor{azulmed}{HTML}{2874A6}
\definecolor{naranja}{HTML}{CA6F1E}
\definecolor{verde}  {HTML}{1E8449}
\definecolor{lila}   {HTML}{6C3483}
\definecolor{rojo}   {HTML}{922B21}
\definecolor{gris}   {HTML}{5D6D7E}
\definecolor{linea}  {HTML}{AEB6BF}
\definecolor{faz}    {HTML}{EAF2F8}
\definecolor{fna}    {HTML}{FEF9E7}
\definecolor{fve}    {HTML}{E9F7EF}
\definecolor{fli}    {HTML}{F4ECF7}
\definecolor{fro}    {HTML}{FDEDEC}
\usepackage{tikz}
\usepackage{pgfplots}
\pgfplotsset{compat=1.18}

\usetikzlibrary{
  babel,
  arrows.meta,calc,patterns,decorations.pathmorphing,
  decorations.markings,shapes.geometric,positioning,
  backgrounds,fit,3d
}
\pgfplotsset{compat=1.18}

\usetikzlibrary{arrows.meta,calc,patterns,decorations.pathmorphing,
  decorations.markings,shapes.geometric,positioning,backgrounds,fit,3d}
\usepackage[many,breakable]{tcolorbox}

\usepackage{verbatim}  

% Definir el entorno de definición

\theoremstyle{definition} % Texto en letra normal (no cursiva)
\newtheorem{definicion}{Definición}[section] % Numeradas por secció
\newtheorem{teorema}{Teorema} [section] % Numerados por sección
\newtheorem{lema}{Lema}[teorema]
\newtheorem{corolario}{Corolario}[teorema]
\newtheorem{ejemplo}{Ejemplo}[subsection]
\newtheorem{proposicion}{Proposición}[section]
\newcommand{\bolaa}[2][\varepsilon]{B_{#1}(#2)}
\newcommand{\bolac}[2][\varepsilon]{{B}_{#1}[#2]}
\usepackage{geometry}
\newcommand{\clau}[1]{\overline{#1}}
\newcommand{\inn}[1]{#1 \in \mathbb{N}}
\newcommand{\kas}[1]{#1 \xrightarrow{k \to \infty}}
\newcommand{\tendsto}[3]{#1 \xrightarrow{#2 \to #3}}
\newcommand{\col}[1]{\begin{pmatrix} #1_1 \\ \vdots \\ #1_n \end{pmatrix}}
\newcommand{\R}{\mathbb{R}}

\newcommand{\N}{\mathbb{N}}
\newcommand{\Q}{\mathbb{Q}}
\newcommand{\norm}[1]{\left\lVert#1\right\rVert}
\newcommand{\abs}[1]{\left\lvert#1\right\rvert}
\newcommand{\inner}[2]{\left\langle#1,\,#2\right\rangle}
\newcommand{\Int}{\operatorname{int}}
\newcommand{\Gr}{\operatorname{Gr}}
\newcommand{\Dom}{\operatorname{Dom}}
\newcommand{\To}{\longrightarrow}
\newcommand{\Mto}{\longmapsto}














\geometry{
  a4paper,
  left=1.5cm,
  right=1.5cm,
  top=2cm,
  bottom=2cm,
  headheight=1cm,
  footskip=1cm
}



\begin{document}  


\newpage
\section{Fórmulas Timer}
\large
\[\alpha = 2^n - \dfrac{T_{\text{overflow}} \cdot F_{\text{osc}}}{4 \cdot \operatorname{prescaler}}\] 
\normalsize
\begin{itemize}
  

\item $\alpha$ es el valor desde el que empieza a contar el timer

\item $n = $ 8/16 bits 

\item $T_{\text{overflow}} =$ tiempo que tarda en hacer overflow (tiempo que queremos temporizar en segundos)

 \item $F_{\text{osc}} = $ frecuencia de oscilación del chip (normalmente $8 \cdot 10^6$Hz) 

\item $\operatorname{prescaler} = $ lo que divide la frecuencia. Tienen que ser potencias de $2$ hasta $128$ 


\end{itemize}
\begin{center}
\begin{tikzpicture}
  % Segmento base
  \draw[thick] (0,0) -- (10,0);

  % Marcas en 0 y 1
  \draw (0,0.15) -- (0,-0.15) node[below] {$0$};
  \draw (10,0.15) -- (10,-0.15) node[below] {$1$};

  % Posición de alpha (ajustable)
  \def\apos{4}

  % Flecha en alpha
  \draw[->, thick] (\apos,1) -- (\apos,0.1);
  \node[above] at (\apos,1) {$\alpha$};

  % Bracket debajo del tramo (alpha,1)
  \draw[decorate,decoration={brace,mirror,amplitude=6pt}]
    (\apos,-0.2) -- (10,-0.2)
    node[midway,below=8pt] {$T_{\text{overflow}}$};

\end{tikzpicture}

\end{center}
Esta fórmula se calcula desde el menor prescaler hasta el mayor. Debemos empezar con $n = 8$. Si sigue dando negativo, probamos con $n = 16$. 
Hay que obtener el primer valor positivo. La idea es que el timer empiece a contar desde $\alpha$ y se produzca un overflow.
\\
En esta práctica nos interesa el $\boldsymbol{\alpha}$ y el $\mathbf{prescaler}$.

\section{Fórmulas AD}
\subsection{Fórmula $\lambda$}
\large
\[ \lambda = \frac{v_{\text{ref+}}-v_{\text{ref-}}}{2^n-1} \]
\normalsize
\begin{itemize}
  \item $\lambda$ es la resolución por la que debemos multiplicar el valor inicial de la lectura de los registros \textbf{ADRESH} y \textbf{ADRESL}
  \item  $n$ es el número de bits que podemos guardar en los registros  \textbf{ADRESH} y \textbf{ADRESL} (en nuestro caso $n = 10$)
  \item $v_{\text{ref+}}$ es el voltaje alto 
  \item $v_{\text{ref-}}$ es el voltaje bajo
\end{itemize}
\subsection{Fórmula $\mathbf{prescaler}$}
\[ \frac{f_{\text{osc}}}{f_{\text{AD}}} = k \]
De este valor $k$, calculamos la siguiente potencia de 2 (máxima hasta 64).
\begin{itemize}
  \item $f_{\text{osc}}$ es la frecuencia del chip (normalmente $8 \cdot 10^6$)
  \item $f_{\text{AD}}$ es la frecuencia interna del AD (normalmente $625 \cdot 10^3$)
\end{itemize}
\begin{ejemplo} Para frecuencia básica del chip:
  \[ \frac{8 \cdot 10^6}{625 \cdot 10^3} = 12.8\]
  Así que ponemos $\operatorname{prescaler} = 16$
\end{ejemplo}
En esta práctica nos interesa $\boldsymbol{\lambda}$ y el valor del $\mathbf{prescaler}$
\end{document}


